% Źródła: Eves s. 284--285, 298; Wikipedia: Giovanni Girolamo Saccheri, Tommaso Ceva, Saccheri quadrilateral.
% https://en.wikipedia.org/wiki/Giovanni_Girolamo_Saccheri

\section{Giovanni Girolamo Saccheri} % lang-pl
\section{Giovanni Girolamo Saccheri} % lang-it
\label{sekcja_saccheri}

Pierwszą osobą, która na serio zapyta, co by było, gdyby piąty postulat okazał się fałszywy, będzie jezuita z Pawii. % lang-pl
W tej sekcji poznamy jego życie, nauczyciela, jego plan i czworokąt, od którego zaczną się wszystkie późniejsze rozważania. % lang-pl
La prima persona a chiedersi seriamente che cosa accadrebbe se il quinto postulato fosse falso sarà un gesuita di Pavia. % lang-it
In questa sezione conosceremo la sua vita, il suo maestro, il suo piano e il quadrilatero da cui prenderanno le mosse tutte le considerazioni successive. % lang-it

Giovanni Girolamo Saccheri urodzi się 5 września 1667 roku w Sanremo, w 1685 roku wstąpi do nowicjatu jezuitów, a święcenia kapłańskie przyjmie w 1694 roku. % lang-pl
Będzie wykładał filozofię w Turynie (1694--1697), a potem filozofię, teologię i matematykę na uniwersytecie w Pawii (1697--1733). % lang-pl
Zamknie oczy 25 października 1733 roku w Mediolanie. % lang-pl
Giovanni Girolamo Saccheri nascerà il 5 settembre 1667 a Sanremo, nel 1685 entrerà nel noviziato dei gesuiti e riceverà l'ordinazione sacerdotale nel 1694. % lang-it
Insegnerà filosofia a Torino (1694--1697) e poi filosofia, teologia e matematica all'Università di Pavia (1697--1733). % lang-it
Chiuderà gli occhi il 25 ottobre 1733 a Milano. % lang-it
\index[persons]{Saccheri, Giovanni Girolamo}%

Giovanni Girolamo Saccheri będzie uczniem Tommasa Cevy, brata Giovanniego, autora twierdzenia Cevy (w tej książce z numerem \ref{twierdzenie_cevy}). % lang-pl
Giovanni Girolamo Saccheri sarà allievo di Tommaso Ceva, fratello di Giovanni, autore del teorema di Ceva (in questo libro numerato come \ref{twierdzenie_cevy}). % lang-it
Tommaso Ceva (1648--1737), jezuita, przez trzydzieści osiem lat będzie uczył matematyki i retoryki w kolegium Brera w Mediolanie; pod jego kierunkiem Saccheri napisze swoją pierwszą pracę. % lang-pl
To Ceva przedstawi Saccheriego swojemu bratu Giovanniemu, z którym młody jezuita utrzyma bliskie kontakty naukowe. % lang-pl
Tommaso Ceva (1648--1737), gesuita, insegnerà per trentotto anni matematica e retorica nel collegio di Brera a Milano; sotto la sua guida Saccheri scriverà il suo primo lavoro. % lang-it
Sarà Ceva a presentare Saccheri al fratello Giovanni, con il quale il giovane gesuita manterrà stretti rapporti scientifici. % lang-it
\index[persons]{Ceva, Tommaso}%
\index[persons]{Ceva, Giovanni}%

W 1733 roku Giovanni Girolamo Saccheri spróbuje zbadać, co by było gdyby piąty postulat Euklidesa okazał się fałszywy. % lang-pl
Jego prace zapoczątkują systematyczne rozważania nad alternatywami dla geometrii euklidesowej. % lang-pl
Nel 1733 Giovanni Girolamo Saccheri cercherà di capire che cosa accadrebbe se il quinto postulato di Euclide fosse falso. % lang-it
I suoi lavori daranno inizio allo studio sistematico delle alternative alla geometria euclidea. % lang-it

Książka nosi tytuł \emph{Euclides ab omni naevo vindicatus} (Euklides oczyszczony ze wszelkiej skazy) i ukaże się tuż przed śmiercią autora. % lang-pl
Saccheri zachwyci się w swojej wcześniejszej książce o logice metodą dowodu nie wprost i zechce zastosować ją do postulatu równoległości \cite[s. 285]{eves1_1972}. % lang-pl
Il libro s'intitolerà \emph{Euclides ab omni naevo vindicatus} (Euclide liberato da ogni macchia) e uscirà poco prima della morte dell'autore. % lang-it
Saccheri, affascinato in un'opera precedente di logica dal metodo della dimostrazione per assurdo, vorrà applicarlo al postulato delle parallele \cite[p. 285]{eves1_1972}. % lang-it

\begin{definition}
[czworokąt Saccheriego] % lang-pl
[quadrilatero di Saccheri] % lang-it
	Czworokąt $ABCD$, w którym kąty $A$ i $B$ są proste, a boki $AD$ i $BC$ są równe. % lang-pl
	Odcinek $AB$ nazywamy podstawą, $DC$ górną podstawą, a kąty $D$ i $C$ kątami przy górnej podstawie. % lang-pl
	Un quadrilatero $ABCD$ in cui gli angoli $A$ e $B$ sono retti e i lati $AD$ e $BC$ sono uguali. % lang-it
	Il segmento $AB$ si chiama base, $DC$ base superiore, e gli angoli $D$ e $C$ angoli alla base superiore. % lang-it
\index{czworokąt!Saccheriego} % lang-pl
\index{quadrilatero!di Saccheri} % lang-it
\end{definition}

Eves \cite[s. 298]{eves1_1972}. % lang-pl
Eves \cite[p. 298]{eves1_1972}. % lang-it

Podobny czworokąt rozważą już wcześniej Sabit ibn Kurra (IX wiek), Omar Chajjam (XI wiek) i, tuż przed Saccherim, Giordano Vitale. % lang-pl
Una figura simile era stata considerata già da Thabit ibn Qurra (IX secolo), Omar Khayyam (XI secolo) e, poco prima di Saccheri, Giordano Vitale. % lang-it

Bez piątego postulatu łatwo pokazać, że kąty przy górnej podstawie są równe; są więc trzy możliwości: oba ostre, oba proste albo oba rozwarte. % lang-pl
Saccheri nazwie je \emph{hipotezą kąta ostrego}, \emph{prostego} i \emph{rozwartego}. % lang-pl
Senza il quinto postulato si mostra facilmente che gli angoli alla base superiore sono uguali; le possibilità sono dunque tre: entrambi acuti, entrambi retti o entrambi ottusi. % lang-it
Saccheri le chiamerà \emph{ipotesi dell'angolo acuto}, \emph{retto} e \emph{ottuso}. % lang-it
\index{hipoteza!kąta ostrego (prostego, rozwartego)} % lang-pl
\index{ipotesi!dell'angolo acuto (retto, ottuso)} % lang-it

Plan jest prosty: pokazać, że hipoteza kąta ostrego i hipoteza kąta rozwartego prowadzą do sprzeczności, więc musi zachodzić hipoteza kąta prostego, a z niej wynika piąty postulat. % lang-pl
Hipotezę kąta rozwartego Saccheri odrzuci szybko, milcząco zakładając, że prosta jest nieskończona -- tak jak zrobił to Euklides. % lang-pl
Z hipotezą kąta ostrego pójdzie znacznie trudniej. % lang-pl
Il piano è semplice: mostrare che l'ipotesi dell'angolo acuto e quella dell'angolo ottuso portano a contraddizione, sicché deve valere l'ipotesi dell'angolo retto, dalla quale discende il quinto postulato. % lang-it
L'ipotesi dell'angolo ottuso Saccheri la scarterà presto, assumendo tacitamente che la retta sia infinita -- come aveva fatto Euclide. % lang-it
Con l'ipotesi dell'angolo acuto andrà molto peggio. % lang-it

Saccheri wyprowadzi z niej długi szereg twierdzeń, które dziś czytamy jako twierdzenia geometrii hiperbolicznej, po czym na siłę wmówi sobie sprzeczność, odwołując się do mglistych pojęć o elementach nieskończonych; wniosek uzna za „sprzeczny z naturą prostej''. % lang-pl
Gdyby przyznał, że sprzeczności nie znalazł, dziś uchodziłby za odkrywcę geometrii nieeuklidesowej. % lang-pl
Wygląda na to, że książkę wkrótce po wydaniu wycofano ze sprzedaży, więc jej wpływ na współczesnych będzie znikomy \cite[s. 285]{eves1_1972}. % lang-pl
Da quest'ipotesi dedurrà una lunga serie di teoremi che oggi leggiamo come teoremi della geometria iperbolica, per poi convincersi a forza di aver trovato una contraddizione, appellandosi a nozioni vaghe su elementi infiniti; la conclusione gli parrà «contraria alla natura della retta». % lang-it
Se avesse ammesso di non aver trovato alcuna contraddizione, oggi sarebbe considerato lo scopritore della geometria non euclidea. % lang-it
Sembra che il libro sia stato ritirato dal commercio poco dopo la pubblicazione, così che il suo influsso sui contemporanei sarà minimo \cite[p. 285]{eves1_1972}. % lang-it

Dopiero w 1889 roku Eugenio Beltrami przypomni tę książkę światu (zobacz sekcję \ref{sekcja_beltrami_cayley_klein}). % lang-pl
Solo nel 1889 Eugenio Beltrami riporterà questo libro all'attenzione del mondo (si veda la sezione \ref{sekcja_beltrami_cayley_klein}). % lang-it
